%!TeX root=../thesis.tex

\chapter{Research Instruments}
\label{apx:instruments}

% TODO: include the actual instruments used in the empirical study.
% This appendix lets the reader see exactly what participants received.

\section{Questionnaire}
\label{apx:questionnaire}

% TODO: paste the full questionnaire here. Example structure:
%
% The following questionnaire was administered to team members after
% the usage period. It was delivered via [Google Forms / LimeSurvey / etc.].
%
% \subsection{Section 1: Demographics}
% \begin{enumerate}
%   \item What is your role in the team? (Developer / Tech Lead / Architect / Other)
%   \item How many years of professional software development experience do you have?
%   \item How familiar were you with software architecture concepts before this study?
%         (5-point Likert: Not at all -- Extremely familiar)
% \end{enumerate}
%
% \subsection{Section 2: Hypo-Stage Usability}
% \begin{enumerate}[resume]
%   \item I found Hypo-Stage easy to use.
%         (5-point Likert: Strongly disagree -- Strongly agree)
%   \item The process of creating a hypothesis in Hypo-Stage was intuitive.
%         (5-point Likert)
%   \item Assessing uncertainty and impact levels was straightforward.
%         (5-point Likert)
% \end{enumerate}
%
% \subsection{Section 3: Open-Ended Questions}
% \begin{enumerate}[resume]
%   \item What did you find most useful about Hypo-Stage?
%   \item What was the most difficult aspect of using Hypo-Stage?
%   \item What improvements would you suggest?
% \end{enumerate}

\section{Interview Guide}
\label{apx:interview-guide}

% TODO: paste the semi-structured interview guide here. Example:
%
% The following topics were used as a guide for semi-structured interviews.
% The interviewer adapted follow-up questions based on participant responses.
%
% \begin{enumerate}
%   \item \textbf{Context:} Can you describe your team's typical approach to
%         architectural decisions before using Hypo-Stage?
%   \item \textbf{Usage:} Walk me through how you used Hypo-Stage in your
%         last sprint/cycle. What hypotheses did you create?
%   \item \textbf{Value:} Did Hypo-Stage change how your team discusses
%         architecture? In what way?
%   \item \textbf{Barriers:} What made it difficult to adopt Hypo-Stage
%         or ArchHypo concepts?
%   \item \textbf{Refinements:} If you could change one thing about the tool,
%         what would it be?
% \end{enumerate}

\section{Informed Consent Form}
\label{apx:consent}

% TODO: include the informed consent form presented to participants.
% This is required by ethics committees and demonstrates research rigor.

\section{Training Communication}
\label{apx:training-communication}

The following message was sent to teams before the training session that introduced
the ArchHypo experiment and the Hypo-Stage plugin in Backstage.

\begin{quote}
\textbf{New Plugin Experiment in Backstage: Hypo Stage --- Practicing the ArchHypo Framework}

Deciding when to make (or delay) a software architecture decision is one of the
toughest daily challenges in engineering, regardless of the project:

\begin{itemize}
  \item Deciding too early can lock the system into the wrong path.
  \item Deciding too late can lead to costly rework and technical debt.
\end{itemize}

In the middle of this lies something that often remains implicit: uncertainty. We
deal with requirements that are not yet clear or solutions we are not sure will
scale or integrate well. This uncertainty is rarely made explicit, and decisions
often end up being built on unvalidated premises.

To evolve past this dilemma, I am proposing the application of a framework to help
us identify, evaluate, and manage architectural uncertainties in a structured way.
We will put this into practice using the Hypo Stage plugin, which I helped develop
as part of an open-source project. It is already available in our Backstage (in
production and under a feature flag for our experiment).

This plugin serves as a tool to practice the ArchHypo (Architectural Hypotheses)
framework: a lightweight approach that applies hypothesis-driven engineering to
represent uncertainties impacting software architecture. It helps us plan actions
to reduce uncertainty or impact before "locking in" a decision.

\textbf{What ArchHypo proposes, in short:}

\textbf{Architectural Hypotheses:} Falsifiable statements expressing team beliefs
about requirements (e.g., quality attributes) or current/candidate solutions. The
focus is on uncertainties that affect architectural decisions, not just business
hypotheses.

\textbf{Evaluation:} Each hypothesis is assessed based on its uncertainty level
(how far we are from proving it true/false) and impact (the consequences and cost
of being wrong). This helps prioritize and choose treatment tactics.

\textbf{Technical Plan:} For each hypothesis, the team defines actions to either
reduce uncertainty (e.g., spikes, experiments, architectural triggers, metrics) or
reduce impact (e.g., modularity, flexibility, isolation of the affected part),
allowing us to delay decisions more safely.

\textbf{Iterative Cycle:} Identify, evaluate, plan, execute, and reassess. The
process is continuous; as new uncertainties emerge, plans are adjusted over time.

The Hypo Stage plugin in Backstage allows you to register hypotheses, link them to
quality attributes, record uncertainty/impact assessments, and maintain technical
plans with dates and progress tracking --- all integrated into our catalog and
workflow.

\textbf{Want to better understand the problem, the framework, and how to use the plugin?}

This Week Day, Calendar Date, at HH:MM during the Meeting Agenda, I will be leading
a session on this topic. We will discuss why finding the right moment for
architectural decisions is so difficult and how ArchHypo can help.

I am looking forward to seeing you there to discuss this experiment and how to use
Hypo Stage in our daily routine.
\end{quote}
